% Figure: upstream-to-downstream process flowsheet for a generic
% biomanufacturing run.
\begin{figure}[htbp]
\centering
\definecolor{engorange}{RGB}{217, 119, 6}
\begin{tikzpicture}[
  block/.style={draw=engblack, thick, rounded corners,
                fill=engblue!10, minimum width=2.3cm, minimum height=0.9cm,
                font=\small, align=center},
  upstream/.style={block, fill=englime!15},
  recovery/.style={block, fill=engorange!15},
  purif/.style={block, fill=engblue!15},
  final/.style={block, fill=engred!10},
  arrow/.style={-{Stealth}, thick},
  lab/.style={font=\scriptsize, engblack!70},
]

% Upstream
\node[upstream] (cellbank) at (0, 0) {Cell bank\\(WCB)};
\node[upstream] (seed) at (3.0, 0) {Seed\\culture};
\node[upstream] (inoc) at (6.0, 0) {Inoculum\\$10\text{--}100\,\text{L}$};
\node[upstream] (production) at (9.0, 0) {Production\\bioreactor};

\draw[arrow] (cellbank) -- (seed);
\draw[arrow] (seed) -- (inoc);
\draw[arrow] (inoc) -- (production);

% Down to recovery
\draw[arrow] (production.south) -- ++(0, -0.7) -| (1.5, -1.6);

% Recovery row
\node[recovery] (clar) at (1.5, -2.2) {Clarification\\(centrifuge / depth filter)};
\node[recovery] (cap) at (5.0, -2.2) {Capture\\(IEX or Protein A)};
\node[recovery] (interm) at (8.5, -2.2) {Intermediate\\polishing};

\draw[arrow] (clar) -- (cap);
\draw[arrow] (cap) -- (interm);

% Down to purification / final
\draw[arrow] (interm.south) -- ++(0, -0.7) -| (1.5, -3.8);

% Purification row
\node[purif] (viralclear) at (1.5, -4.4) {Viral\\clearance};
\node[purif] (polish) at (4.5, -4.4) {Final\\polishing};
\node[purif] (uf) at (7.5, -4.4) {Ultrafiltration /\\buffer exchange};

\draw[arrow] (viralclear) -- (polish);
\draw[arrow] (polish) -- (uf);

% Final
\draw[arrow] (uf.south) -- ++(0, -0.7) -| (3.0, -6.0);

\node[final] (formulation) at (3.0, -6.4) {Formulation};
\node[final] (filtration) at (6.0, -6.4) {Sterile\\filtration};
\node[final] (fill) at (9.0, -6.4) {Fill /\\finish};

\draw[arrow] (formulation) -- (filtration);
\draw[arrow] (filtration) -- (fill);

% Labels
\node[lab, anchor=west] at (-0.5, 0.8) {\textbf{Upstream}};
\node[lab, anchor=west] at (-0.5, -1.4) {\textbf{Recovery}};
\node[lab, anchor=west] at (-0.5, -3.6) {\textbf{Purification}};
\node[lab, anchor=west] at (-0.5, -5.6) {\textbf{Final}};

% Yield arrow on the side
\draw[engred, very thick, ->] (11.0, -0.4) -- (11.0, -6.6);
\node[lab, engred, rotate=-90, anchor=south] at (11.3, -3.5) {cumulative yield drops};

\end{tikzpicture}
\caption{Generic upstream-to-finish biomanufacturing process flowsheet.
The upstream train (green) takes the working cell bank through seed
expansion to a production bioreactor whose volume depends on the
product and the cell platform. Recovery (orange) clarifies the broth
and runs a capture chromatography that concentrates the product and
removes the majority of process-related impurities. Purification
(blue) adds viral clearance, polishing, and buffer exchange. The final
train (pink) formulates the bulk drug substance, sterile-filters, and
fills the finished container. Each step has a yield less than unity,
so cumulative recovery from broth to vial multiplies and falls
substantially across the train; the red arrow on the right marks the
direction. The flowsheet is generic, but no commercial process omits
more than one of these compartments.}
\label{fig:vol06ch10:flowsheet}
\end{figure}
